\documentclass{article}
\usepackage[utf8]{inputenc}
\usepackage{geometry}
\usepackage{graphicx}
\usepackage{pgfplots}
\usepackage{amsmath}
\usepackage{amssymb}
\usepackage{amsfonts}
\usepackage[thmmarks, amsmath, thref]{ntheorem}
\usepackage{mdframed}
\usepackage{amsthm}
\usepackage{physics}
\usepackage{proof}
\usepackage{derivative}
\usepackage{hyperref}
\usepackage{wrapfig}
\usepackage{amsthm}
\usepackage{xcolor}
\renewcommand{\theenumi}{\roman{enumi}}

\theoremstyle{nonumberplain}
\theoremheaderfont{\itshape}
\theorembodyfont{\upshape}
\theoremseparator{.}
\theoremsymbol{\ensuremath{\square}}
\theoremsymbol{\ensuremath{\blacksquare}}
\newtheorem{solution}{Solution}
\theoremseparator{. ---}
\theoremsymbol{\mbox{\texttt{;o)}}}

\pgfplotsset{compat=1.18}
\begin{document}

(a) Mostrar que la integral de línea de la función escalar \( f(x, y)\left(f: \mathbb{R}^{2} \longrightarrow \mathbb{R}\right) \) a lo largo de una trayectoria \( \Gamma \subset \mathbb{R}^{2} \) dada en coordenadas polares por \( r=r(\theta) \) con \( \theta_{1} \leq \theta \leq \theta_{2} \), es:
\[
\int_{\theta_{1}}^{\theta_{2}} f(r \cos \theta, r \operatorname{sen} \theta)\left(r^{2}+\left(\frac{\mathrm{d} r}{\mathrm{~d} \theta}\right)^{2}\right)^{\frac{1}{2}} \mathrm{~d} \theta
\]

En coordenadas cartesianas, la curva se escribe como,
\[
  x(\theta) = r(\theta) \cos \theta, \quad y(\theta) = r(\theta) \sin \theta
  \]

La integral de línea de \( f \) a lo largo de \( \Gamma \) es:
\[
\int_\Gamma f(x, y) \, ds
\]
donde \( ds \) es el elemento diferencial de longitud de arco. Notemos que el elemento diferencial de arco es:
\[
ds = \sqrt{\left(\frac{dx}{d\theta}\right)^2 + \left(\frac{dy}{d\theta}\right)^2} \, d\theta
\]
Calculamos las derivadas:
\[
\frac{dx}{d\theta} = \frac{d}{d\theta} \left[ r(\theta) \cos \theta \right] = r'(\theta) \cos \theta - r(\theta) \sin \theta
\]
\[
\frac{dy}{d\theta} = \frac{d}{d\theta} \left[ r(\theta) \sin \theta \right] = r'(\theta) \sin \theta + r(\theta) \cos \theta
\]
Sumamos los cuadrados:
\[
\left(\frac{dx}{d\theta}\right)^2 + \left(\frac{dy}{d\theta}\right)^2 = \left[ r'(\theta) \cos \theta - r(\theta) \sin \theta \right]^2 + \left[ r'(\theta) \sin \theta + r(\theta) \cos \theta \right]^2
\]
Expandimos:
\[
= \left(r'(\theta)\right)^2 \cos^2 \theta - 2 r'(\theta) r(\theta) \cos \theta \sin \theta + r^2(\theta) \sin^2 \theta
\]
\[
+ \left(r'(\theta)\right)^2 \sin^2 \theta + 2 r'(\theta) r(\theta) \sin \theta \cos \theta + r^2(\theta) \cos^2 \theta
\]
Agrupamos términos:\\

\( -2 r'(\theta) r(\theta) \cos \theta \sin \theta \) y \( +2 r'(\theta) r(\theta) \sin \theta \cos \theta \) se cancelan.\\
\( (r'(\theta))^2 (\cos^2 \theta + \sin^2 \theta) = (r'(\theta))^2 \)\\
\( r^2(\theta) (\sin^2 \theta + \cos^2 \theta) = r^2(\theta) \)\\

De aquí que...
\[
\left(\frac{dx}{d\theta}\right)^2 + \left(\frac{dy}{d\theta}\right)^2 = (r'(\theta))^2 + r^2(\theta)
\]
Por lo tanto,
\[
ds = \sqrt{r^2(\theta) + \left(r'(\theta)\right)^2} \, d\theta
\]
Vemos la expresión de la integral de línea,
\[
\int_\Gamma f(x, y) \, ds = \int_{\theta_1}^{\theta_2} f(r(\theta) \cos \theta, r(\theta) \sin \theta) \sqrt{r^2(\theta) + (r'(\theta))^2} \, d\theta
\]
Esto es justamente lo que se debía demostrar:
\[
\int_{\theta_{1}}^{\theta_{2}} f(r \cos \theta, r \operatorname{sen} \theta)\left(r^{2}+\left(\frac{\mathrm{d} r}{\mathrm{~d} \theta}\right)^{2}\right)^{\frac{1}{2}} \mathrm{~d} \theta
\]


(b) Calcule la longitud de la cardioide \( r(\theta)=1+\cos \theta \operatorname{con} 0 \leq \theta \leq 2 \pi \).

La longitud \( L \) de una curva en polares es,
\[
L = \int_{\theta_1}^{\theta_2} \sqrt{r^2(\theta) + (r'(\theta))^2} \, d\theta
\]
Para la cardioide:

\( r(\theta) = 1 + \cos \theta \)\\
\( r'(\theta) = -\sin \theta \)

Así sustituyendo:
\[
L = \int_{0}^{2\pi} \sqrt{(1+\cos \theta)^2 + (\sin \theta)^2} \, d\theta
\]
Desarrollamos el integrando:
\[
(1+\cos \theta)^2 + (\sin \theta)^2 = 1 + 2\cos \theta + \cos^2 \theta + \sin^2 \theta
\]
Pero \( \cos^2 \theta + \sin^2 \theta = 1 \), así que:
\[
= 1 + 2\cos \theta + 1 = 2 + 2\cos \theta
\]
Entonces:
\[
L = \int_{0}^{2\pi} \sqrt{2 + 2\cos \theta} \, d\theta
\]
Factorizamos:
\[
\sqrt{2 + 2\cos \theta} = \sqrt{2} \sqrt{1 + \cos \theta}
\]
Y usando la identidad trigonométrica:
\[
1 + \cos \theta = 2 \cos^2 \left(\frac{\theta}{2}\right)
\]
Así que:
\[
\sqrt{1 + \cos \theta} = \sqrt{2} \left|\cos \left(\frac{\theta}{2}\right)\right|
\]
Para \( 0 \leq \theta \leq 2\pi \), el valor absoluto se puede manejar dividiendo el intervalo, pero la función es simétrica.
Por lo tanto:
\[
L
=\int_{0}^{2\pi}2\bigl|\cos\!\tfrac\theta2\bigr|\,d\theta
=4\int_{0}^{\pi}\bigl|\cos u\bigr|\,du
\quad(u=\tfrac\theta2,\;d\theta=2\,du)
\]
\[
=4\Bigl[\int_{0}^{\pi/2}\cos u\,du
-\int_{\pi/2}^{\pi}\cos u\,du\Bigr]
=4\bigl[1-0+(0-(-1))\bigr]
=8.
\]


La longitud de la cardioide \( r(\theta) = 1 + \cos \theta \) para \( 0 \leq \theta \leq 2\pi \) es 8.

\end{document}



